\section{Introduction}

Selenium is incorporated into proteins primarily as selenocysteine (Sec), the 21st amino acid, through recoding of the uracil--guanine--adenine (UGA) stop codon.\cite{Berry1991UGA,Kryukov2003Selenoproteomes,Shetty2015Selenocysteine,hondal_selenium_2011,Labunskyy2014} Relative to cysteine (Cys), Sec has a lower side-chain pK$_a$, a larger and more polarizable chalcogen atom, and a selenolate form that is typically more nucleophilic than the corresponding thiolate, which together make selenium especially effective in redox catalysis.\cite{johansson_selenocysteine_2005,hondal_selenium_2011,Arner2010,ChungKrahn2022SelenocysteineToolbox,Reich2016,Kanavos2026SecSelectivity} The glutathione peroxidases (GPXs) are a major family of thiol-dependent peroxidases that reduce \ch{H2O2} and organic hydroperoxides to water or the corresponding alcohols, and they include some of the earliest recognized vertebrate selenoenzymes.\cite{Rotruck1973,BrigeliusFlohe2013GlutathionePeroxidases,pei_research_2023} Glutathione peroxidase~6 (GPX6) is particularly unusual within this family because the catalytic residue differs across mammalian lineages. Sec is retained in humans and some other primate lineages, whereas Cys--containing GPX6 has evolved convergently across mammals, including rodents, rabbits, cats, and some monkeys.\cite{BrigeliusFlohe2013GlutathionePeroxidases,rees_ancient_2024}

Mechanistic studies of seleno-GPXs by Morokuma and coworkers\cite{prabhakar_elucidation_2005,prabhakar_is_2006} indicate that peroxide reduction proceeds predominantly through a stepwise pathway in which the resting selenol is first deprotonated to the reactive selenolate, followed by O--O bond cleavage and formation of selenenic acid. 
In this framework, formation of the selenolate carries a substantial free-energy barrier and precedes peroxide bond cleavage, whereas the alternative concerted pathway, in which Se--H and O--O bond cleavage occur in a single water-assisted step, was calculated to have a substantially higher barrier.\cite{prabhakar_elucidation_2005,prabhakar_is_2006}
Because this acid--base step is the earliest chemically discriminating event between Sec-- and Cys--containing enzymes, it provides a focused point of comparison across GPX6 variants.\cite{prabhakar_is_2006,cardey_selenocysteine_2007} The analogous thiol-to-thiolate transformation can be defined for Cys--containing homologs, allowing a consistent mechanistic framework across the evolutionary series.\cite{cardey_selenocysteine_2007}
The stepwise reaction scheme considered here is summarized in Figure~\ref{fig:gpx6_mechanism}.


\begin{figure}[htbp]
\centering
\includegraphics[width=\textwidth]{Fig1.png}
\caption[Stepwise mechanism for hydrogen peroxide reduction]{Modeled I$\rightarrow$II$\rightarrow$III pathway for \ch{H2O2} reduction showing peptide-bound Sec49 and Gln83. \textbf{I}, Sec49 selenol with neutral Gln83 and \ch{H2O2}; \textbf{II}, Sec49 selenolate with O-protonated Gln83 and \ch{H2O2}; \textbf{III}, selenenic acid, water, and regenerated neutral Gln83. TS~I--II and TS~II--III denote the corresponding transition-state regions. Curved arrows show nucleophilic attack by peroxide oxygen on the Sec49 proton and by Gln83 O$\varepsilon$ on the peroxide proton in TS~I--II, and attack by Se$^{-}$ on peroxide oxygen and by the second peroxide oxygen on the Gln83 proton in TS~II--III.}
\label{fig:gpx6_mechanism}
\end{figure}

Previous phylogenetic and functional analyses identified a clear evolutionary split in GPX6. Human GPX6 (UniProt Knowledgebase [UniProtKB] P59796) retains Sec at the catalytic site, whereas mouse GPX6 (UniProtKB Q91WR8) belongs to a rodent Cys--containing lineage.\cite{rees_ancient_2024} Ancestral sequence reconstruction further suggests that the last common ancestor of these lineages encoded Sec and that Cys--containing GPX6 variants arose through multiple independent Sec--loss events across mammals.\cite{rees_ancient_2024} These convergent GPX6-Cys changes show signatures of positive selection, suggesting that evolution favored the Cys state and that the Cys--containing enzyme may have acquired or improved a still-unresolved novel function.\cite{rees_ancient_2024} Reintroduction of Sec into rodent GPX6 restores classical peroxidase activity, underscoring the catalytic importance of the reactive residue.\cite{rees_ancient_2024} However, a simple reversal from Cys to Sec is genetically demanding because Sec incorporation requires not only a UGA codon at the catalytic position but also the selenocysteine insertion machinery, including a selenocysteine insertion sequence (SECIS) ribonucleic acid (RNA) element in the 3$'$ untranslated region.\cite{Berry1991UGA,Kryukov2003Selenoproteomes,Shetty2015Selenocysteine,ChungKrahn2022SelenocysteineToolbox} GPX6 thus allows us to examine the same catalytic substitution in two sequence backgrounds that have followed different evolutionary histories. The unresolved question is how the surrounding residues alter the effect of the Sec--to-Cys exchange and whether low-barrier routes connect the corresponding catalytic states.

This background dependence is a form of epistasis: the effect of a mutation depends on the other residues present in the protein. It can restrict the evolutionary routes available and make later changes depend on earlier ones.\cite{maynardsmith1970proteinspace,Weinreich2006,starr2016epistasis,storz_compensatory_2018} Experimental and structural studies have related this path dependence to conformational and biophysical interactions among substitutions.\cite{Ortlund2007,Bridgham2009,Harms2014} In enzymes, substitutions can alter electrostatics, hydrogen-bonding networks, conformational preferences, and pK$_a$ values, and thereby change transition-state stabilization and catalytic barriers.\cite{Warshel2006,Kamerlin2010,VillaWarshel2001,Tokuriki2009,storz_compensatory_2018} The effect of replacing Sec with Cys in GPX6 may therefore depend on the surrounding sequence rather than on the reactive residue alone.\cite{rees_ancient_2024,starr2016epistasis}

In this work, we use empirical valence bond (EVB) simulations to quantify how sequence divergence and epistatic interactions reshape the catalytic free-energy landscape connecting human and mouse GPX6. We first evaluate the proton-transfer step that generates the reactive selenolate or thiolate species in Sec-- and Cys--containing variants of both homologs. We then examine how individual and combinatorial substitutions modulate activation free energies, use these energetic profiles to construct low-barrier mutational pathways, and place the resulting trajectories in the phylogenetic context established previously.\cite{rees_ancient_2024} This design tests whether the Sec/Cys transition behaves as a simple catalytic-residue exchange or as a background-dependent reorganization of transition-state stabilization. As a secondary analysis, we also ask whether sequence generation from EVB-derived structures recovers ancestral-like sequence features more effectively than the energetic pathways alone.
The active-site architecture is shown in Figure~\ref{fig:gpx6_structure_mutations}; the full structural view used to map the 20 proximal human/mouse substitutions is provided in Figure~\ref{SI-fig:gpx6_structure_mutations_full}.

\begin{figure}[htbp]
\centering
\includegraphics[width=0.88\textwidth]{Fig2.png}
\caption[Structural comparison of the GPX6 active site]{Detail of the superposed GPX6 active sites. The experimental mouse Cys-GPX6 structure (orange) is from Protein Data Bank (PDB) entry~7FC2 (UniProtKB Q91WR8), and the human Sec-GPX6 structure (white) is an AlphaFold~2 model of UniProtKB P59796. The two residues directly included in the modeled proton transfer, the catalytic chalcogen residue and Gln83, are shown together with neighboring Trp157 and peroxide (PRX).}
\label{fig:gpx6_structure_mutations}
\end{figure}
